\scanpage{124}
\margin{\textcircled{30}}
\uline{Généralités sur Lefschetz}. \uline{Lefschetz-Verdier}

\uline{Préliminaires à la formule de \deleted{\unsure{Künneth}} Lefschetz}

Classique, on a une variété compacte $X$, un endomorphisme $f$, d'où (si on prend la cohomologie à coeff. dans un \uline{corps} de base $K$) des $f^{(i)}$ des $H^i(X)$, et la formule dit
\[
\sum(-1)^i\operatorname{Tr}f^{(i)}=\text{nb de pts fixes de }f
\]
avec grain de sel pour le dernier mb (si p.ex. les pts fixes sont isolés, il faut leur attribuer une multiplicité de façon convenable).
\margin{\unsure{Notons}}

En fait, on peut souvent interpréter le $2^{\text{ième}}$ mb \inserted{et prouver la formule} $\unsure{=}$ si $X$ n'est pas une variété (p.ex. $X$ un complexe simplicial fini), et Verdier a généralisé \inserted{par une méthode et un \unsure{langage} très \unsure{différents}} au cas où on a un faisceau de coefficients $\pm$ quelconque (soumis à des \unsure{seules|simples} conditions de nature ``\unsure{finitude}''). On a maintenant une \unsure{dém.} très simple dans cas ``élémentaire'' (variété, coeff constants, ou du moins loc. constants), comme résultant \deleted{de la} essentiellement du formalisme de la dualité de Poincaré de façon classique.

\inserted{Dans le contexte des \unsure{schémas topologiques|schémas algébriques}} On donnera l'approche de Verdier \inserted{(utilisant des moyens plus \unsure{profds})} sans entrer dans les détails, et nous \deleted{\unsure{nous}} donnerons en détail le ``cas élémentaire'', qui nous permettra ensuite
\mspage{123}
